\chapter{Regresja liniowa: estymacja, diagnostyka i predykcja}
\label{ch:linear-regression}
\paragraph{Powiązania i odsyłacze.} Metody predykcji i doboru złożoności rozwija rozdział~\ref{ch:model-selection}, a teorię testów~\ref{ch:parametric-tests}.


\section{Kiedy statystyka średniej przestaje wystarczać?}
Porównując trzy technologie, w rozdziale~\ref{ch:parametric-tests}
analizowaliśmy różnice średnich między grupami. Co jednak zrobić, gdy
sprawność urządzenia zależy także od temperatury, czasu pracy i wieku
maszyny? Potrzebujemy modelu opisującego \emph{warunkową średnią}
wyniku przy ustalonych cechach. Regresja liniowa jest najprostszą
konstrukcją o tej własności. Jej nazwa odnosi się do liniowości
\emph{współczynników}, niekoniecznie do liniowości wszystkich funkcji
wykorzystanych jako zmienne objaśniające: można używać $x^2$ lub
interakcji $xz$, a nadal otrzymać model liniowy względem parametrów.

\begin{definicja}[Model liniowy przy ustalonym planie doświadczenia]
\label{def:linear-fixed}
Niech $n,p\in\N$, $n>p\ge1$, a
$X\in\R^{n\times p}$ będzie znaną macierzą o rzędzie $p$.
Niech $\beta\in\R^p$, zaś $\varepsilon\in\R^n$ będzie
losowym wektorem o $\E\varepsilon=0$ i
$\Cov(\varepsilon)=\sigma^2I_n$, gdzie $\sigma^2>0$.
W \emph{liniowym modelu homoskedastycznym} obserwujemy
$Y=X\beta+\varepsilon$. Jeśli dodatkowo
$\varepsilon\sim\mathcal N_n(0,\sigma^2I_n)$, mówimy o
\emph{normalnym modelu liniowym}. Macierz $X$ to
\emph{macierz planu}, a $\beta$ to wektor współczynników.
Wiersz $x_i^\top$ macierzy $X$ opisuje cechy $i$-tej jednostki.
\end{definicja}

Przy losowych cechach $X$ wszystkie poniższe twierdzenia można
stosować \emph{warunkowo względem $X$}, o ile odpowiednie założenia
spełnione są prawie na pewno. Sam fakt losowości cech nie unieważnia
regresji. Potrzebujemy jednak warunkowych założeń
$\E(\varepsilon\mid X)=0$ i
$\Cov(\varepsilon\mid X)=\sigma^2I_n$,
aby bez zmian stosować twierdzenia o nieobciążoności
i wariancji warunkowej. Same warunki bezwarunkowe
$\E\varepsilon=0$ i $\Cov(\varepsilon)=\sigma^2I_n$
nie wystarczają. Przykładowo dla skalarnych
$X\sim\Normal(0,1)$ i $\varepsilon=X$ mamy
$\E\varepsilon=0$, ale $\E(\varepsilon\mid X)=X$.
Problem ten występuje m.in. przy pominięciu
istotnej zmiennej zakłócającej.

\begin{przyklad}[Co oznacza współczynnik w modelu?]
Jeśli $\E(Y_i\mid X)=\beta_0+\beta_1t_i+\beta_2g_i$,
gdzie $t_i$ to temperatura, a $g_i\in\{0,1\}$ oznacza
technologię B, to $\beta_1$ jest zmianą średniej odpowiedzi
na wzrost temperatury o jedną jednostkę \emph{przy tej samej
technologii}, a $\beta_2$ jest różnicą średnich B--A
\emph{przy tej samej temperaturze}. Gdy znaczenie temperatury
może zależeć od technologii, sprawdzamy rozszerzenie
$\beta_3t_ig_i$. Wtedy różnica B--A to
$\beta_2+\beta_3t_i$ i nie wolno przypisywać
$\beta_2$ znaczenia stałej różnicy między grupami.
To zapowiada analizę kowariancji (ANCOVA).
\end{przyklad}

\section{Metoda najmniejszych kwadratów jako rzut ortogonalny}
\begin{twierdzenie}[Istnienie, jednoznaczność i postać MNK]
\label{thm:ols-formula}
Załóżmy, że $X\in\R^{n\times p}$ ma rząd $p\le n$.
Dla dowolnego $y\in\R^n$ funkcja
$Q_y(b)=\|y-Xb\|_2^2$ ma dokładnie jedno minimum
na $\R^p$, osiągane w punkcie
\[
 \widehat\beta=(X^\top X)^{-1}X^\top y.
\]
Przy $H=X(X^\top X)^{-1}X^\top$ wektor
$\widehat y=Hy$ jest rzutem ortogonalnym $y$ na
$\operatorname{im}X$, a reszta $e=(I_n-H)y$ jest
prostopadła do każdej kolumny $X$.
\end{twierdzenie}
\begin{proof}
Macierz $X^\top X$ jest dodatnio określona, ponieważ
dla $v\ne0$ zachodzi
$v^\top X^\top Xv=\|Xv\|_2^2>0$.
Zdefiniujmy $b_*=(X^\top X)^{-1}X^\top y$.
Wtedy $X^\top(y-Xb_*)=0$, więc $y-Xb_*$ jest
prostopadły do wszystkich wektorów postaci $Xv$.
Dla dowolnego $b$ z twierdzenia Pitagorasa otrzymujemy
$\|y-Xb\|_2^2=\|y-Xb_*\|_2^2+\|X(b-b_*)\|_2^2$.
Drugi składnik jest zerowy tylko dla $b=b_*$.
Ponadto $H^\top=H$ i $H^2=H$, zatem $H$ jest
rzutem ortogonalnym na $\operatorname{im}X$.
\end{proof}

\begin{uwaga}[Dlaczego nie obliczamy odwrotności w programie?]
Wzór teoretyczny zawiera $(X^\top X)^{-1}$, ale
jawne wyznaczanie tej odwrotności jest niepotrzebne i
bywa niestabilne numerycznie, ponieważ liczba uwarunkowania
$X^\top X$ jest kwadratem liczby uwarunkowania $X$.
W praktyce stosujemy rozkład QR lub SVD (rozkład według
wartości osobliwych). W przykładzie programistycznym
użyjemy \texttt{numpy.linalg.lstsq}. Analogicznie,
przy regularizacji ridge w rozdziale~\ref{ch:model-selection}
nie będziemy formować odwrotności macierzy wprost.
\end{uwaga}

\section{Dlaczego współczynniki MNK bywają optymalne?}
\begin{twierdzenie}[Nieobciążoność i macierz kowariancji]
\label{thm:ols-cov}
W modelu~\ref{def:linear-fixed} przy pełnym rzędzie $X$
estymator $\widehat\beta=(X^\top X)^{-1}X^\top Y$
spełnia
\[
 \E\widehat\beta=\beta,\qquad
 \Cov(\widehat\beta)=\sigma^2(X^\top X)^{-1}.
\]
\end{twierdzenie}
\begin{proof}
Z tożsamości $\widehat\beta-\beta=A\varepsilon$ dla
$A=(X^\top X)^{-1}X^\top$ wynika
$\E\widehat\beta-\beta=A\E\varepsilon=0$.
Ponadto $\Cov(A\varepsilon)=A(\sigma^2I_n)A^\top
=\sigma^2(X^\top X)^{-1}$, ponieważ
$AA^\top=(X^\top X)^{-1}X^\top X(X^\top X)^{-1}$.
\end{proof}

\begin{twierdzenie}[Gauss--Markow: najlepszy liniowy nieobciążony]
\label{thm:gauss-markov}
W założeniach twierdzenia~\ref{thm:ols-cov}, niech
$\widetilde\beta=CY$ dla ustalonej macierzy
$C\in\R^{p\times n}$ takiej, że $CX=I_p$.
Wtedy różnica
$\Cov(\widetilde\beta)-\Cov(\widehat\beta)$
jest nieujemnie określona. Zatem MNK ma najmniejszą
macierz kowariancji w porządku Loewnera \emph{w klasie
liniowych nieobciążonych estymatorów}.
Do twierdzenia nie jest potrzebna normalność błędów.
\end{twierdzenie}
\begin{proof}
Niech $A=(X^\top X)^{-1}X^\top$ i $D=C-A$.
Z równości $CX=AX=I_p$ wynika $DX=0$.
Ponieważ $AD^\top=(X^\top X)^{-1}X^\top D^\top
=(X^\top X)^{-1}(DX)^\top=0$, mamy
$CC^\top=AA^\top+DD^\top$.
W konsekwencji
$\Cov(CY)-\Cov(AY)=\sigma^2DD^\top$,
a dla każdego $v\in\R^p$ zachodzi
$v^\top DD^\top v=\|D^\top v\|_2^2\ge0$.
\end{proof}

\begin{uwaga}[Czego twierdzenie Gaussa--Markowa nie dowodzi?]
Nie oznacza ono, że MNK jest najlepszy spośród wszystkich
estymatorów, odporny na dowolne odstające obserwacje ani że
współczynniki mają rozkład normalny. To ostatnie wymaga
założenia normalności (lub odpowiedniej asymptotyki).
Jeśli $\Cov(\varepsilon)=\Omega\ne\sigma^2I_n$,
to klasyczna minimalność MNK nie wynika z twierdzenia;
po właściwym opisaniu $\Omega$ można rozważyć ważone
najmniejsze kwadraty.
\end{uwaga}

\section{Reszty, wariancja i dokładne rozkłady w modelu normalnym}
\begin{twierdzenie}[Rozkład współczynników i sumy kwadratów reszt]
\label{thm:ols-gaussian}
Niech $n>p\ge1$, $X$ ma pełny rząd kolumnowy i
$Y\sim\mathcal N_n(X\beta,\sigma^2I_n)$, $\sigma>0$.
Niech $e=(I_n-H)Y$, $\mathrm{RSS}=\|e\|_2^2$ oraz
$s^2=\mathrm{RSS}/(n-p)$. Wtedy:
\begin{enumerate}
\item $\widehat\beta\sim\mathcal N_p\!\bigl(\beta,
 \sigma^2(X^\top X)^{-1}\bigr)$;
\item $\mathrm{RSS}/\sigma^2\sim\chi^2_{n-p}$;
\item $\widehat\beta$ i $\mathrm{RSS}$ są niezależne,
 a $\E(s^2)=\sigma^2$.
\end{enumerate}
\end{twierdzenie}
\begin{proof}
Estymator jest liniowym przekształceniem wektora normalnego,
więc pierwszy punkt wynika z obliczonej już macierzy kowariancji.
Macierz $H$ jest rzutem ortogonalnym rzędu $p$, zatem
$I_n-H$ jest rzutem ortogonalnym rzędu $n-p$.
Niech $U$ będzie macierzą ortogonalną, która diagonalizuje
$I_n-H$ do $\operatorname{diag}(I_{n-p},0_p)$.
Wektor $U^\top\varepsilon/\sigma$ ma niezależne
współrzędne $\mathcal N(0,1)$. Ponieważ
$(I_n-H)X\beta=0$, otrzymujemy
$\mathrm{RSS}/\sigma^2=\sum_{j=1}^{n-p}Z_j^2$.
Niezależność wynika z ortogonalności rzutów:
$\Cov(H\varepsilon,(I_n-H)\varepsilon)=0$,
a wspólna normalność oznacza niezależność wektorów.
Estymator zależy jedynie od $HY$, zaś RSS jedynie od
$(I_n-H)Y$. Wreszcie średnia rozkładu
$\chi^2_{n-p}$ wynosi $n-p$.
\end{proof}

\begin{wniosek}[Dokładny przedział dla współczynnika]
\label{cor:ols-ci}
W założeniach twierdzenia~\ref{thm:ols-gaussian}
niech $c\in\R^p\setminus\{0\}$ będzie ustalonym
wektorem. Wtedy
\[
 T_c=\frac{c^\top\widehat\beta-c^\top\beta}
 {s\sqrt{c^\top(X^\top X)^{-1}c}}\sim t_{n-p}.
\]
W rezultacie przedział
$c^\top\widehat\beta\pm t_{n-p,1-\alpha/2}
 s\sqrt{c^\top(X^\top X)^{-1}c}$
ma pokrycie $1-\alpha$, dla $0<\alpha<1$.
\end{wniosek}
\begin{proof}
Licznik podzielony przez
$\sigma\sqrt{c^\top(X^\top X)^{-1}c}$ ma rozkład
$\mathcal N(0,1)$; zgodnie z twierdzeniem
\ref{thm:ols-gaussian} jest niezależny od
$(n-p)s^2/\sigma^2\sim\chi^2_{n-p}$.
Ich iloraz ma rozkład $t_{n-p}$.
\end{proof}

\begin{przyklad}[Dlaczego zmiana jednostek nie zmienia istotności?]
Jeżeli temperatura jest mierzona w stopniach Celsjusza,
a następnie zastąpimy ją różnicą skali, np. podzielimy
przez $10$, odpowiedni współczynnik zostanie pomnożony
przez $10$, podobnie jak jego błąd standardowy.
Statystyka t dla testu zerowego współczynnika pozostanie
niezmieniona. Interpretując efekt, zawsze podajemy
\emph{jednostkę} zmiennej objaśniającej.
\end{przyklad}

\section{Predykcja przyszłej obserwacji a szacowanie średniej}
\begin{twierdzenie}[Przedział ufności i przedział predykcji]
\label{thm:prediction-vs-mean}
Przy założeniach twierdzenia~\ref{thm:ols-gaussian}
ustalmy niezerowy wektor cech $x_0\in\R^p\setminus\{0\}$ i niech
$Y_0=x_0^\top\beta+\varepsilon_0$, gdzie
$\varepsilon_0\sim\Normal(0,\sigma^2)$ jest niezależna
od $Y$. Niech $h_0=x_0^\top(X^\top X)^{-1}x_0$.
Wówczas przedział ufności dla średniej warunkowej
$\mu_0=x_0^\top\beta$ wynosi
$x_0^\top\widehat\beta\pm t_{n-p,1-\alpha/2}s\sqrt{h_0}$,
a przedział predykcji dla nowego $Y_0$ ma postać
\[
 x_0^\top\widehat\beta\pm
 t_{n-p,1-\alpha/2}s\sqrt{1+h_0}.
\]
Każdy z przedziałów ma odpowiednio pokrycie
$1-\alpha$ dla własnego celu; drugi dotyczy \emph{losowej
nowej obserwacji}, nie stałej wartości parametru.
\end{twierdzenie}
\begin{proof}
Z poprzedniego twierdzenia
$x_0^\top(\widehat\beta-\beta)$ jest normalny o
wariancji $\sigma^2h_0$. Dla nowego wyniku błąd
$Y_0-x_0^\top\widehat\beta$ jest sumą niezależnych
składników $\varepsilon_0$ i
$-x_0^\top(\widehat\beta-\beta)$, ma więc wariancję
$\sigma^2(1+h_0)$ i rozkład normalny.
Oba błędy są niezależne od $s^2$, bo $\varepsilon_0$
jest niezależne od próby, a estymator i RSS są niezależne.
Standaryzacja przez $s$ prowadzi do rozkładu t.
\end{proof}

\begin{uwaga}[Ekstrapolacja]
Gdy $x_0$ jest daleko od zebranych obserwacji, wartość
$h_0$ może być duża; dodatkowo liniowa zależność może
w tym obszarze nie obowiązywać. Nawet węższy przedział
ufności dla średniej nie gwarantuje prawidłowości modelu
poza zakresem danych. \emph{Ekstrapolacja} wymaga
uzasadnienia dziedzinowego, nie tylko małej wartości p.
\end{uwaga}

\section{Ogólny test F i dlaczego ANOVA jest regresją}
\begin{twierdzenie}[Test zagnieżdżonych modeli normalnych]
\label{thm:partial-f-test}
Niech $Y\sim\mathcal N_n(\mu,\sigma^2I_n)$,
$\sigma>0$, a $V_0\subsetneq V_1$ będą ustalonymi
podprzestrzeniami $\R^n$ wymiarów odpowiednio
$p_0<p_1<n$. Oznaczmy ich rzuty ortogonalne przez
$H_0,H_1$ i resztowe sumy kwadratów przez
$\mathrm{RSS}_j=\|(I_n-H_j)Y\|_2^2$.
Przy hipotezie zerowej $\mu\in V_0$ statystyka
\[
 F=\frac{(\mathrm{RSS}_0-\mathrm{RSS}_1)/(p_1-p_0)}
 {\mathrm{RSS}_1/(n-p_1)}\sim F_{p_1-p_0,n-p_1}.
\]
\end{twierdzenie}
\begin{proof}
Ponieważ $V_0\subset V_1$, macierz $H_1-H_0$ jest
rzutem ortogonalnym na $V_1\cap V_0^\perp$ rzędu
$p_1-p_0$. Macierz $I_n-H_1$ jest rzutem na
$V_1^\perp$ rzędu $n-p_1$. Są to rzuty ortogonalne
o iloczynie zero; wektor średniej $\mu\in V_0$ leży
w jądrach obu. Z twierdzenia o ortogonalnych składowych
wektora normalnego kwadraty norm przeskalowane przez
$\sigma^2$ są niezależnymi zmiennymi o rozkładach
$\chi^2_{p_1-p_0}$ i $\chi^2_{n-p_1}$.
Ponieważ $\mathrm{RSS}_0-\mathrm{RSS}_1
=\|(H_1-H_0)Y\|_2^2$, ich iloraz ma podany rozkład.
\end{proof}

\begin{przyklad}[Jednoczynnikowa ANOVA jako test F regresji]
Niech $G_i\in\{1,\ldots,k\}$ oznacza grupę obserwacji
i załóżmy, że w każdej grupie jest przynajmniej jedna
jednostka. W modelu ograniczonym $V_0$ wszystkie grupy
mają wspólną średnią; w pełnym $V_1$ każda grupa ma
osobną średnią. Wymiary tych przestrzeni wynoszą
odpowiednio $1$ i $k$. Resztowe sumy kwadratów to
$\mathrm{RSS}_0=\sum_i(Y_i-\bar Y)^2$ i
$\mathrm{RSS}_1=\sum_g\sum_{i:G_i=g}(Y_i-\bar Y_g)^2$.
Różnica jest sumą kwadratów między grupami,
$\sum_gn_g(\bar Y_g-\bar Y)^2$.
Podstawienie do twierdzenia~\ref{thm:partial-f-test}
odtwarza statystykę ANOVA z rozdziału~\ref{ch:parametric-tests}.
\end{przyklad}

\begin{uwaga}[ANCOVA jest kolejnym krokiem]
Jeśli do macierzy planu dodamy kolumnę temperatury,
test ogólnej różnicy grup kontrolując temperaturę jest
porównaniem zagnieżdżonych przestrzeni kolumn. Warto
sprawdzić interakcję grupy z temperaturą przed przyjęciem
jednego wspólnego nachylenia. Nie należy utożsamiać
istotności współczynnika z przyczynowym efektem
zmiany technologii bez argumentu o randomizacji
lub identyfikacji przyczynowej.
\end{uwaga}

\section{Diagnostyka modelu: co może zepsuć poprawne wzory?}
\begin{definicja}[Dźwignia obserwacji i reszta studencka]
Dla macierzy planu $X$ o pełnym rzędzie i rzutowania
$H=X(X^\top X)^{-1}X^\top$ liczbę $h_{ii}=H_{ii}$
nazywamy \emph{dźwignią} $i$-tej obserwacji.
Przy $h_{ii}<1$ i oszacowaniu $s>0$ definiujemy
\emph{wewnętrznie standaryzowaną resztę}
$r_i=e_i/(s\sqrt{1-h_{ii}})$.
\end{definicja}

\begin{stwierdzenie}[Dlaczego obserwacje skrajne mogą mieć małe reszty?]
W modelu $Y=X\beta+\varepsilon$ z
$\Cov(\varepsilon)=\sigma^2I_n$ reszta $e=(I_n-H)Y$
spełnia $\E e=0$ i
$\Cov(e)=\sigma^2(I_n-H)$.
W szczególności $\Var(e_i)=\sigma^2(1-h_{ii})$,
a $0\le h_{ii}\le1$ oraz $\sum_ih_{ii}=p$.
\end{stwierdzenie}
\begin{proof}
Równość $(I_n-H)X=0$ zapewnia
$e=(I_n-H)\varepsilon$. Korzystając z
$(I_n-H)^2=I_n-H=(I_n-H)^\top$,
otrzymujemy macierz kowariancji. Z własności rzutu
$h_{ii}=e_i^\top He_i=\|He_i\|_2^2\in[0,1]$,
a jego ślad jest równy rzędowi $p$.
\end{proof}

\paragraph{Pięć pytań diagnostycznych.}
\begin{enumerate}
\item \emph{Czy średnia jest właściwie opisana?}
Rysujemy reszty względem wartości przewidywanych i cech,
sprawdzając krzywiznę i interakcje. To kontrola poprawności
postaci modelu, a nie tylko test pojedynczego współczynnika.
\item \emph{Czy wariancja jest stała?} Jeśli rozrzut rośnie
wraz z poziomem sygnału, klasyczne błędy standardowe mogą być
niewłaściwe. Rozważamy ważone MNK, odporne oszacowania
kowariancji lub inny model rozkładu.
\item \emph{Czy obserwacje są niezależne?} Przy powtarzanych
pomiarach osoby lub silnej autokorelacji potrzebna może być
kowariancja klastrowa albo jawny model zależności.
\item \emph{Czy istnieją punkty wpływowe?} Analizujemy
\emph{dźwignie}, reszty i zmianę współczynników po usunięciu
obserwacji; sam duży $|r_i|$ nie wyczerpuje diagnostyki.
\item \emph{Czy kolumny planu są rozróżnialne?}
Silna współliniowość może znacznie powiększyć
$\operatorname{diag}((X^\top X)^{-1})$, a przy dokładnej
zależności liniowej $\widehat\beta$ nie jest jednoznaczny.
\end{enumerate}

\begin{stwierdzenie}[Korekta odporna na heteroskedastyczność]
\label{prop:hetero-ols}
Jeśli $Y=X\beta+\varepsilon$, $\E\varepsilon=0$ oraz
$\Cov(\varepsilon)=\Omega$ dla znanej nieujemnie
określonej macierzy $\Omega$, to
\[
 \Cov(\widehat\beta)=(X^\top X)^{-1}X^\top
 \Omega X(X^\top X)^{-1}.
\]
Gdy błędy są nieskorelowane, lecz mają różne wariancje,
można \emph{oszacowywać} $\Omega$ diagonalnie poprzez
$e_i^2/(1-h_{ii})^2$ dla $h_{ii}<1$. Otrzymana
macierz kanapkowa HC3 jest narzędziem asymptotycznym;
nie daje automatycznie dokładnych małopróbkowych testów t.
\end{stwierdzenie}
\begin{proof}
Wciąż $\widehat\beta-\beta
=(X^\top X)^{-1}X^\top\varepsilon$.
Wzór na kowariancję liniowego przekształcenia daje tezę.
Zastąpienie nieznanych elementów $\Omega$ oszacowaniami
nie jest tożsamością skończonopróbkową i wymaga osobnej
analizy asymptotycznej.
\end{proof}

\begin{przyklad}[Pominięta zmienna i pozorna istotność]
Niech prawdziwa zależność warunkowej średniej ma postać
$\E(Y\mid x,z)=\beta_0+\beta_1x+\gamma z$.
Jeżeli regresujemy $Y$ wyłącznie na $1$ i $x$,
to średnia estymowanego nachylenia, przy ustalonych
wektorach $x,z$ i błędach o średniej zero, wynosi
\[
 \E\widehat b_1=\beta_1+
 \gamma\frac{\sum_i(x_i-\bar x)(z_i-\bar z)}
                   {\sum_i(x_i-\bar x)^2}.
\]
Współczynnik może być istotnie różny od zera nawet
gdy $\beta_1=0$, jeśli $z$ jest związane z $x$.
Sama większa próba nie usuwa tego rodzaju obciążenia.
\end{przyklad}

\section{Studium praktyczne: jak czytać wynik regresji?}
Rozważmy dane z pomiarów procesu technologicznego:
czas działania w godzinach $t_i=1,\ldots,8$ i wyniki
$y=(4{,}8,5{,}1,6{,}0,6{,}1,7{,}2,7{,}4,8{,}1,8{,}9)$.
Modelujemy $\E(Y_i\mid t_i)=\beta_0+\beta_1t_i$.
Przed interpretacją parametru $\beta_1$ pytamy, czy
każdy pomiar dotyczy innej, niezależnej partii, czy tej
samej maszyny śledzonej w czasie. W drugim przypadku
standardowa analiza niezależnych błędów może nie być
poprawna, choć dopasowana prosta będzie wyglądała tak samo.

\begin{lstlisting}[language=Python,caption={Kompletny program: MNK, przedział dla nachylenia, diagnostyka dźwigni i predykcja (NumPy, SciPy).}]
import numpy as np
from scipy.stats import t as student

t = np.arange(1.0, 9.0)
y = np.array([4.8, 5.1, 6.0, 6.1, 7.2, 7.4, 8.1, 8.9])
X = np.column_stack((np.ones_like(t), t))
n, p = X.shape

beta, _, rank, _ = np.linalg.lstsq(X, y, rcond=None)
assert rank == p
residual = y - X @ beta
rss = residual @ residual
s2 = rss / (n-p)
xtx_inv = np.linalg.solve(X.T @ X, np.eye(p))
se = np.sqrt(np.diag(s2 * xtx_inv))
quant = student.ppf(0.975, n-p)
ci_slope = (beta[1] - quant*se[1], beta[1] + quant*se[1])

H_diag = np.sum((X @ xtx_inv) * X, axis=1)
standardized_residual = residual / np.sqrt(s2 * (1-H_diag))
x0 = np.array([1.0, 9.0])
mean_prediction = x0 @ beta
h0 = x0 @ xtx_inv @ x0
pi = (mean_prediction - quant*np.sqrt(s2*(1+h0)),
      mean_prediction + quant*np.sqrt(s2*(1+h0)))
print('coefficients:', beta)
print('95% CI slope:', ci_slope)
print('leverages:', H_diag)
print('standardized residuals:', standardized_residual)
print('prediction interval at t=9:', pi)
\end{lstlisting}

\begin{lstlisting}[language=Julia,caption={MNK poprzez rozkład QR i błąd standardowy w Julii (standardowe LinearAlgebra, Statistics).}]
using LinearAlgebra, Statistics

t = collect(1.0:8.0)
y = [4.8, 5.1, 6.0, 6.1, 7.2, 7.4, 8.1, 8.9]
X = hcat(ones(length(t)), t)
beta = X \ y  # stabilniejsze niz inv(X'X) * X'y
r = y - X * beta
n, p = size(X)
s2 = sum(abs2, r) / (n - p)
xtx_inverse = (X' * X) \ Matrix{Float64}(I, p, p)
se = sqrt.(diag(s2 * xtx_inverse))
println((beta=beta, se=se, rss=sum(abs2, r)))
\end{lstlisting}

\paragraph{Interpretacja programu.}
Dla podanych danych program otrzymuje \(\widehat\beta_0\approx4{,}075\)
i \(\widehat\beta_1\approx0{,}5833\), a przedział ufności
$95\%$ dla nachylenia to w przybliżeniu
$[0{,}5035;\ 0{,}6632]$. W modelu zakładającym
niezależne normalne błędy przedział predykcji nowego
wyniku dla $t_0=9$ wynosi około $[8{,}669;\ 9{,}981]$.
Są to \emph{współczynniki modelu warunkowej średniej},
nie gotowe twierdzenia przyczynowe. Przedział predykcji zawiera składnik
wariancji nowej obserwacji i jest zwykle szerszy
niż przedział ufności dla średniej. Jeżeli reszty
wykazują autokorelację lub heteroskedastyczność,
nie należy raportować klasycznych wartości p bez
ponownego oszacowania niepewności.

\section{Wybór metody i kontrola jakości analizy}
\begin{center}\small
\begin{tabular}{>{\raggedright\arraybackslash}p{.29\textwidth}>{\raggedright\arraybackslash}p{.31\textwidth}>{\raggedright\arraybackslash}p{.29\textwidth}}
\toprule
\textbf{Sytuacja} & \textbf{Rozważana metoda} & \textbf{Ograniczenia}\\
\midrule
Wynik ciągły, średnia liniowa, błędy o stałej wariancji & Zwykłe MNK; przy normalności dokładne t i F & Zależność, nieliniowość średniej, punkty wpływowe.\\
Zmienna wynikowa ciągła, niejednakowe wariancje & MNK z odpornym SE lub WLS & WLS wymaga uzasadnienia wag; HC3 jest przybliżeniem.\\
Silna zależność w czasie albo grupowanie pomiarów & Model korelacji, GLS, modele mieszane lub odporny SE klastrowy & Liczba niezależnych jednostek i specyfikacja zależności.\\
Wpływ skrajnych wyników & Regresja odporna, kontrola wpływu, analiza wrażliwości & Nie usuwać punktów wyłącznie na podstawie p-value.\\
Dane binarne lub zliczenia & GLM z odpowiednią rodziną & Przedział wartości średniej, wariancja zależna od średniej.\\
\bottomrule
\end{tabular}
\end{center}

\begin{zadanie}[Czy dwa modele mówią to samo?]
W badaniu stwierdzono, że współczynnik temperatury jest
istotny w regresji bez zmiennej opisującej wiek maszyny,
a nieistotny po jej dodaniu. Czy to dowód, że temperatura
nie wpływa na wynik?
\par\emph{Rozwiązanie.} Nie. W pierwszym modelu współczynnik
może odzwierciedlać również skorelowany z temperaturą wiek.
W drugim rośnie liczba współczynników i może wzrosnąć
błąd standardowy. Właściwe jest porównanie założeń,
przedziałów ufności, rozkładu cech oraz celu analizy,
a przy twierdzeniach przyczynowych także projektu badania.
Sama zmiana przekroczenia progu $p=0{,}05$ nie dowodzi
braku efektu.
\end{zadanie}

% ================================================================
% ROZDZIAL 14 -- GLM
% ================================================================
